\chapter{Making Decisions in Python: if/else Statements}

\section{Introduction: Real-Life Hook}

In daily life, students make decisions every day:  

\begin{itemize}
    \item Should I bring an umbrella? (If it rains)  
    \item Should I do homework now or later?  
    \item Should I eat a snack before lunch?  
\end{itemize}

Programs can also make decisions using \textbf{conditions}.

\section{Concept: Conditional Statements}

\begin{definition}{Conditional Statement}
A \textbf{conditional statement} allows a program to make decisions based on whether a condition is true or false.
\end{definition}

\subsection{Python if Statement}

The simplest form:

\begin{verbatim}
age = 15
if age >= 18:
    print("You can vote")
\end{verbatim}

- If the condition (\texttt{age >= 18}) is true, the program executes the indented block.  
- If false, nothing happens.

\subsection{Python if/else Statement}

To handle both cases:

\begin{verbatim}
age = 15
if age >= 18:
    print("You can vote")
else:
    print("You cannot vote yet")
\end{verbatim}

Output:

\begin{verbatim}
You cannot vote yet
\end{verbatim}

\section{Python if/elif/else for Multiple Choices}

\begin{verbatim}
marks = 75
if marks >= 90:
    print("Grade: A+")
elif marks >= 80:
    print("Grade: A")
elif marks >= 70:
    print("Grade: B+")
else:
    print("Grade: C or below")
\end{verbatim}

\subsection{Real-Life Analogy: Deciding What to Wear}

\begin{itemize}
    \item If it is raining $\rightarrow$ wear raincoat  
    \item Else if it is cold $\rightarrow$ wear sweater  
    \item Else $\rightarrow$ wear T-shirt  
\end{itemize}

\section{Step-by-Step Diagram: Decision Making}

\begin{center}
\begin{tikzpicture}[node distance=2cm]
\node (start) [draw, ellipse, fill=green!20] {Start Program};
\node (condition) [draw, diamond, fill=yellow!20, below of=start] {Is condition true?};
\node (trueblock) [draw, rectangle, fill=orange!20, below left of=condition, xshift=-2cm] {Execute if True};
\node (falseblock) [draw, rectangle, fill=orange!20, below right of=condition, xshift=2cm] {Execute if False};
\node (end) [draw, ellipse, fill=green!20, below of=condition, yshift=-4cm] {End Program};
\draw[->] (start) -- (condition);
\draw[->] (condition.south west) -- (trueblock.north);
\draw[->] (condition.south east) -- (falseblock.north);
\draw[->] (trueblock.south east) -- (end.north);
\draw[->] (falseblock.south west) -- (end.north);
\end{tikzpicture}
\end{center}

\section{Mini-Project: Exam Result Checker}

Create a Python program that asks the user for their marks and prints if they passed or failed:

\begin{verbatim}
marks = int(input("Enter your marks: "))
if marks >= 40:
    print("You passed the exam!")
else:
    print("You failed the exam.")
\end{verbatim}

\section{Key Concepts}

\begin{itemize}
    \item Conditional statements allow programs to make decisions.
    \item Use \texttt{if}, \texttt{else}, and \texttt{elif} to handle different conditions.
    \item Indentation is important in Python to define blocks of code.
    \item Decision-making makes programs dynamic and interactive.
\end{itemize}

\section{End-of-Chapter Exercises}

\begin{exercise}{Very Short Questions}
\begin{enumerate}
    \item What is a conditional statement?
    \item Which Python keyword is used for multiple choices?
    \item What happens if an if-condition is false and there is no else block?
\end{enumerate}
\end{exercise}

\begin{exercise}{Short Answer Questions}
\begin{enumerate}
    \item Write a Python program that prints "Even" if a number is even, otherwise "Odd".
    \item Explain the difference between if/else and if/elif/else.
    \item Why is indentation important in Python?
\end{enumerate}
\end{exercise}

\begin{exercise}{Activity-Based Questions}
\begin{enumerate}
    \item Create a program that asks the user for the weather and gives a clothing suggestion.
    \item Modify the exam result program to also print grades (A+, A, B+, etc.) using if/elif/else.
    \item Test your program with different inputs to check all branches of decisions.
\end{enumerate}
\end{exercise}

\section{Teacher Notes}

\marginnote{
\textbf{Learning Objectives:} 
\begin{itemize}
    \item Students learn to use if, else, and elif statements.
    \item Students understand decision-making in programs.
    \item Students can create interactive Python programs with multiple outcomes.
\end{itemize}

\textbf{Teaching Tips:} 
\begin{itemize}
    \item Use real-life decision examples in class.
    \item Demonstrate programs with different inputs to show branching.
\end{itemize}

\textbf{Assessment:} 
\begin{itemize}
    \item Check correctness of decision-making programs.
    \item Observe student understanding of conditions and branching.
\end{itemize}
}